APR 01 - CODIGO PIPELINE
import numpy as np
import pandas as pd
df=pd.read_csv("diabetes.csv")
X=df[["Glucose","BMI","Age"]].copy()
T=df["Outcome"].copy()
rng=np.random.default_rng(42)
ids=rng.permutation(len(X))
n=int(.7*len(X))
itr,ite=ids[:n],ids[n:]
Xtr=X.iloc[itr].copy()
Xte=X.iloc[ite].copy()
Ttr=T.iloc[itr].to_numpy()
Tte=T.iloc[ite].to_numpy()
for c in ["Glucose","BMI"]:
    med=Xtr.loc[Xtr[c]!=0,c].median()
    Xtr.loc[Xtr[c]==0,c]=med
    Xte.loc[Xte[c]==0,c]=med
for c in Xtr.columns:
    lo,hi=Xtr[c].min(),Xtr[c].max()
    if hi!=lo:
        Xtr[c]=(Xtr[c]-lo)/(hi-lo)
        Xte[c]=(Xte[c]-lo)/(hi-lo)
A=np.c_[np.ones(len(Xtr)),Xtr.to_numpy()]
W=np.linalg.lstsq(A,Ttr,rcond=None)[0]
B=np.c_[np.ones(len(Xte)),Xte.to_numpy()]
Y=(B@W>=.5).astype(int)
TP=np.sum((Tte==1)&(Y==1))
TN=np.sum((Tte==0)&(Y==0))
FP=np.sum((Tte==0)&(Y==1))
FN=np.sum((Tte==1)&(Y==0))
acc=(TP+TN)/len(Y)
se=TP/(TP+FN) if TP+FN else 0
sp=TN/(TN+FP) if TN+FP else 0
f1=2*TP/(2*TP+FP+FN) if 2*TP+FP+FN else 0
print(acc,se,sp,f1)
# Pipeline do PDF utiliza equacao normal com inv.
# lstsq evita falha por matriz singular.